% V033 证明 | 保留 proofbox 环境
\begin{proofbox}
\textbf{\textcolor{CProof}{待证结论}}
设 $O=(0,0)$、$A=(x_1,y_1)$、$B=(x_2,y_2)$，证明
\[\boxed{S_{\triangle OAB}=\frac12|x_1y_2-x_2y_1|.}\]
以下先给考场最快法，再并列保留四种\textbf{仅用初中知识}的最低门槛证明，最后给常规方法与平行四边形方法。

\medskip
\textbf{\textcolor{CProof}{方法一：最快证明（向量有向面积，考场两步）}}
\begin{center}
\begin{tikzpicture}[x=1cm,y=1cm,font=\small,>=Stealth,
  every node/.style={inner sep=3pt},line cap=round,line join=round]
\path[use as bounding box] (-1,-1.6) rectangle (11,5.5);
\node[font=\small\bfseries,text=CStatement] at (5,5.1) {交换两条向量：转向相反，普通面积相同};

\foreach \shift/\al/\bl/\sgn/\ac/\bc in {0/A/B/+1/CStatement/CExample,5/B/A/-1/CExample/CStatement}{
\begin{scope}[shift={(\shift+1,0.5)},scale=1.3]
\fill[CStatement!17] (0,0)--(2,0)--(0,2)--cycle;
\draw[->,gray] (-0.3,0)--(2.7,0) node[right] {$x$};
\draw[->,gray] (0,-0.3)--(0,2.6) node[above] {$y$};
\draw[->,very thick,\ac] (0,0)--(2,0) node[below] {$\al$};
\draw[->,very thick,\bc] (0,0)--(0,2) node[left] {$\bl$};
\draw (2,0)--(0,2);
\node[below left] at (0,0) {$O$};
\node[align=center] at (0.65,0.6) {$S=\frac12$};
\node at (1.4,3.1) {$\det(OA,OB)=\sgn$};
\end{scope}}

\node[align=center] at (5,-1.15) {左：$A=(1,0),B=(0,1)$；右：$A=(0,1),B=(1,0)$。};
\end{tikzpicture}
\end{center}

\textbf{适用信号：} 已学平面向量的二阶行列式表示，三角形有顶点在原点。
\begin{enumerate}[leftmargin=*,itemsep=0.25em]
\item 写 $\vec{OA}=(x_1,y_1)$、$\vec{OB}=(x_2,y_2)$。
\item 已知二维向量有向面积恒等式 $[\vec{OA},\vec{OB}]=x_1y_2-x_2y_1$，
故 $2S=|[\vec{OA},\vec{OB}]|$，即得所证公式。
\end{enumerate}
\textbf{说明：} 此法\textbf{调用}已建立的向量面积恒等式，是求题最快法；
不能反过来把“行列式等于面积”当作对该恒等式来源的证明。来源见方法二至七。

\medskip
\end{proofbox}
\clearpage
\begin{proofbox}
\textbf{\textcolor{CProof}{方法二：过 B 作竖直线（初中共底面积法）}}
\begin{center}
\begin{tikzpicture}[x=1cm,y=1cm,font=\small,>=Stealth,
  every node/.style={inner sep=3pt},line cap=round,line join=round]
\path[use as bounding box] (-1,-1.6) rectangle (11,5.5);
\node[font=\small\bfseries,text=CStatement] at (5,5.1) {竖直公共底：两段水平高相加};

\coordinate (O) at (2,0); \coordinate (A) at (5,4);
\coordinate (B) at (3,2); \coordinate (E) at (3,1.333333);
\fill[CStatement!18] (O)--(B)--(E)--cycle;
\fill[CExample!20] (A)--(B)--(E)--cycle;
\draw[thick] (O)--(A)--(B)--cycle;
\draw[very thick,CTrap] (B)--(E);
\draw[dashed] (2,0)--(3,0) (5,4)--(3,4);
\draw[<->,CProof] (2,-0.25)--node[below] {$1$}(3,-0.25);
\draw[<->,CProof] (3,4.35)--node[above] {$2$}(5,4.35);
\draw[dotted] (3,-0.25)--(3,4.35);
\node[below left] at (O) {$O(0,0)$};
\node[right] at (A) {$A(3,4)$};
\node[left] at (B) {$B(1,2)$};
\node[right] at (E) {$E_v(1,\frac43)$};
\node[align=left] at (8,2.4) {$BE_v=\frac23$\\[6pt]水平高：$1+2=3$\\[6pt]$S=\frac12\cdot\frac23\cdot3=1$};

\node[align=center] at (5,-1.15) {$O,E_v,A$ 共线；$BE_v$ 将原三角形分成蓝、绿两片。};
\end{tikzpicture}
\end{center}

\textbf{图形构造：} 过 $B$ 作竖直线交\emph{直线} $OA$ 于 $E_v$。
先取 $0<x_2<x_1$ 且 $y_1\ne0$，此时 $E_v$ 在线段 $OA$ 内。
\begin{enumerate}[leftmargin=*,itemsep=0.25em]
\item \textbf{为什么用相似？} $O,E_v,A$ 共线，横坐标之比等于纵坐标之比：
$y_{E_v}=\frac{x_2}{x_1}y_1$，因而 $BE_v=|y_2-\frac{x_2}{x_1}y_1|$。
\item \textbf{为什么面积相加？} $E_v\in OA$，由竖直公共底边 $BE_v$ 将原三角形分成两个：
\[
 S=\tfrac12 BE_v\cdot x_2+\tfrac12 BE_v\cdot(x_1-x_2)
 =\tfrac12 |x_1|\,BE_v
 =\tfrac12|x_1y_2-x_2y_1|.
\]
\end{enumerate}
\textbf{位置扩展：} $E_v$ 在 $OA$ 延长线上时，改用大三角形减小三角形；
令 $t=x_2/x_1$。若 $t<0$，水平高之差为 $|x_1|(1-t)-|x_1|(-t)=|x_1|$；若 $t>1$，差为 $|x_1|t-|x_1|(t-1)=|x_1|$。$t=0,1$ 时一片退化，仍为 $S=\frac12|x_1|BE_v$。
$ x_1=0$ 时不能除以 $x_1$，而应以竖直 $OA$ 为底，直接得 $S=\tfrac12|y_1x_2|$。

\medskip
\end{proofbox}
\clearpage
\begin{proofbox}
\textbf{\textcolor{CProof}{方法三：过 B 作水平线（初中共底面积法）}}
\begin{center}
\begin{tikzpicture}[x=1cm,y=1cm,font=\small,>=Stealth,
  every node/.style={inner sep=3pt},line cap=round,line join=round]
\path[use as bounding box] (-1,-1.6) rectangle (11,5.5);
\node[font=\small\bfseries,text=CStatement] at (5,5.1) {水平公共底：两段竖直高相加};

\coordinate (O) at (2,0); \coordinate (A) at (5,4);
\coordinate (B) at (3,2); \coordinate (E) at (3.5,2);
\fill[CStatement!18] (O)--(B)--(E)--cycle;
\fill[CExample!20] (A)--(B)--(E)--cycle;
\draw[thick] (O)--(A)--(B)--cycle;
\draw[very thick,CTrap] (B)--(E);
\draw[dashed] (2,0)--(6,0) (3.5,2)--(6,2) (5,4)--(6,4);
\draw[<->,CProof] (5.7,0)--node[right] {$2$}(5.7,2);
\draw[<->,CProof] (5.7,2)--node[right] {$2$}(5.7,4);
\node[below left] at (O) {$O(0,0)$};
\node[above] at (A) {$A(3,4)$};
\node[left] at (B) {$B(1,2)$};
\node[below right] at (E) {$E_h(\frac32,2)$};
\node[align=left] at (8.5,2.4) {$BE_h=\frac12$\\[6pt]竖直高：$2+2=4$\\[6pt]$S=\frac12\cdot\frac12\cdot4=1$};

\node[align=center] at (5,-1.15) {与竖直法使用同一三角形；公共底与高的方向互换。};
\end{tikzpicture}
\end{center}

\textbf{图形构造：} 过 $B$ 作水平线交\emph{直线} $OA$ 于 $E_h$。
先取 $0<y_2<y_1$，此时 $E_h\in OA$。
\begin{enumerate}[leftmargin=*,itemsep=0.25em]
\item \textbf{为什么相似换了方向？} 由直线 OA 上的坐标比例
$x_{E_h}=\frac{y_2}{y_1}x_1$，$BE_h=|x_2-\frac{y_2}{y_1}x_1|$。
\item \textbf{为什么也能合并？} 以水平线段 $BE_h$ 为共同底边，两个竖直高为 $y_2$、$y_1-y_2$：
\[
 S=\tfrac12 BE_h\,y_2+\tfrac12 BE_h(y_1-y_2)
 =\tfrac12 |y_1|\,BE_h
 =\tfrac12|x_2y_1-x_1y_2|.
\]
\end{enumerate}
\textbf{位置扩展：} 令 $t=y_2/y_1$。若 $0\le t\le1$，两竖直高之和为 $|y_1|$；若 $t<0$，两高之差为 $|y_1|[(1-t)-(-t)]$；若 $t>1$，两高之差为 $|y_1|[t-(t-1)]$。三种情况均得 $S=\frac12|y_1|BE_h$。
若 $y_1=0$，$OA$ 已是水平线段，$S=\tfrac12|x_1y_2|$。
\textbf{两种初中法的共同不变量：} 分割或相减后，两个垂直于公共底的高的\emph{有向合量}，恰好是 $OA$ 的水平或竖直分量。

\medskip
\end{proofbox}
\clearpage
\begin{proofbox}
\textbf{\textcolor{CProof}{方法四：过 A 作竖直线（初中共底面积法）}}
\begin{center}
\begin{tikzpicture}[x=1cm,y=1cm,font=\small,>=Stealth,
  every node/.style={inner sep=3pt},line cap=round,line join=round]
\path[use as bounding box] (-1,-1.6) rectangle (11,5.5);
\node[font=\small\bfseries,text=CStatement] at (5,5.1) {过 A 作竖直线：交点在 OB 延长线上，面积相减};
\begin{scope}[shift={(2,0)},scale=0.7]
\coordinate (O) at (0,0); \coordinate (A) at (3,4);
\coordinate (B) at (1,2); \coordinate (F) at (3,6);
\fill[CStatement!18] (O)--(A)--(B)--cycle;
\fill[CTrap!20] (A)--(B)--(F)--cycle;
\draw[thick] (O)--(A)--(B)--cycle;
\draw[dashed,thick] (B)--(F);
\draw[very thick,CTrap] (A)--(F);
\draw[dotted] (3,0)--(A);
\draw[dashed] (O)--(3,0) (B)--(3,2);
\draw[<->,CProof] (0,-0.35)--node[below] {$3$}(3,-0.35);
\draw[<->,CProof] (1,1.7)--node[below] {$2$}(3,1.7);
\node[below left] at (O) {$O(0,0)$};
\node[right] at (A) {$A(3,4)$};
\node[left] at (B) {$B(1,2)$};
\node[above] at (F) {$F_v(3,6)$};
\end{scope}
\node[align=left] at (8,2.5) {$AF_v=2$\\[7pt]水平高之差：$3-2=1$\\[7pt]$S=\frac12\cdot2\cdot(3-2)=1$};
\node[align=center] at (5,-1.15) {$O,B,F_v$ 依次共线；$S_{OAB}=S_{OAF_v}-S_{BAF_v}$。};
\end{tikzpicture}
\end{center}

\textbf{构造与相似：} 当 $x_2\ne0$，过 $A$ 作竖直线交直线 $OB$ 于 $F_v$。
由 $O,B,F_v$ 共线的坐标比例，
\[
 y_{F_v}=\frac{x_1}{x_2}y_2,\qquad
 AF_v=\left|y_1-\frac{x_1y_2}{x_2}\right|.
\]
\textbf{图中的面积差：} $F_v=(3,6)$ 在射线 $OB$ 上且越过 $B$。
以 $AF_v$ 为公共底，大三角形 $OAF_v$ 的水平高为 $3$，小三角形 $BAF_v$ 的水平高为 $2$，故
\[
 S_{OAB}=S_{OAF_v}-S_{BAF_v}
 =\frac12\cdot2\cdot(3-2)=1.
\]
\textbf{一般化：} 令 $t=x_1/x_2$。$0\le t\le1$ 时 $F_v$ 在线段 $OB$ 内，
两片面积相加，水平高之和为 $|x_2|$。
$t<0$ 时水平高之差为 $|x_2|[(1-t)-(-t)]$；
$t>1$ 时为 $|x_2|[t-(t-1)]$，两者也都等于 $|x_2|$。
因此无论交点在内部还是延长线上，
\[
 S=\frac12|x_2|\,AF_v
 =\frac12|x_2y_1-x_1y_2|
 =\boxed{\frac12|x_1y_2-x_2y_1|}.
\]
\textbf{零分母：} 若 $x_2=0$，$OB$ 本身竖直，直接以其为底，
$S=\frac12|y_2x_1|$；若 $B=O$，面积为零。
\end{proofbox}
\clearpage
\begin{proofbox}
\textbf{\textcolor{CProof}{方法五：过 A 作水平线（初中共底面积法）}}
\begin{center}
\begin{tikzpicture}[x=1cm,y=1cm,font=\small,>=Stealth,
  every node/.style={inner sep=3pt},line cap=round,line join=round]
\path[use as bounding box] (-1,-1.6) rectangle (11,5.5);
\node[font=\small\bfseries,text=CStatement] at (5,5.1) {过 A 作水平线：公共底不变，两竖直高相减};
\begin{scope}[shift={(2,0)},scale=0.7]
\coordinate (O) at (0,0); \coordinate (A) at (3,4);
\coordinate (B) at (1,2); \coordinate (F) at (2,4);
\fill[CStatement!18] (O)--(A)--(B)--cycle;
\fill[CTrap!20] (A)--(B)--(F)--cycle;
\draw[thick] (O)--(A)--(B)--cycle;
\draw[dashed,thick] (B)--(F);
\draw[very thick,CTrap] (A)--(F);
\draw[dashed] (O)--(5.3,0) (B)--(4.3,2) (A)--(5.3,4);
\draw[<->,CProof] (5,0)--node[right] {$4$}(5,4);
\draw[<->,CProof] (4,2)--node[right] {$2$}(4,4);
\node[below left] at (O) {$O(0,0)$};
\node[above right] at (A) {$A(3,4)$};
\node[left] at (B) {$B(1,2)$};
\node[above left] at (F) {$F_h(2,4)$};
\end{scope}
\node[align=left] at (8.3,2.5) {$AF_h=1$\\[7pt]竖直高之差：$4-2=2$\\[7pt]$S=\frac12\cdot1\cdot(4-2)=1$};
\node[align=center] at (5,-1.15) {$O,B,F_h$ 依次共线；减去红片 $\triangle BAF_h$，留下蓝片 $\triangle OAB$。};
\end{tikzpicture}
\end{center}

\textbf{构造与相似：} 当 $y_2\ne0$，过 $A$ 作水平线交直线 $OB$ 于 $F_h$。
利用纵坐标之比，
\[
 x_{F_h}=\frac{y_1}{y_2}x_2,\qquad
 AF_h=\left|x_1-\frac{y_1x_2}{y_2}\right|.
\]
\textbf{图中的面积差：} $F_h=(2,4)$ 在 $OB$ 的延长线上。
以 $AF_h=1$ 为公共底，$O$ 与 $B$ 到水平底的高分别为 $4$、$2$，所以
\[
 S_{OAB}=S_{OAF_h}-S_{BAF_h}
 =\frac12\cdot1\cdot(4-2)=1.
\]
\textbf{一般化：} 令 $t=y_1/y_2$。$0\le t\le1$ 时面积相加，两竖直高之和为 $|y_2|$；
$t<0$ 时两高之差为 $|y_2|[(1-t)-(-t)]$；
$t>1$ 时为 $|y_2|[t-(t-1)]$。三种情形都给出
\[
 S=\frac12|y_2|\,AF_h
 =\frac12|x_1y_2-x_2y_1|.
\]
\textbf{零分母：} 若 $y_2=0$，直接以水平 $OB$ 为底，$S=\frac12|x_2y_1|$；
若 $B=O$，面积为零。
\textbf{四种作法的共同点：} 过 $A$ 或过 $B$、水平或竖直均可；
由相似求交点，再按交点位置加减共底面积，最终得到同一个公式。
\end{proofbox}
\clearpage
\begin{proofbox}
\textbf{\textcolor{CProof}{方法六：常规证明（底边 $OA$ + 点到直线距离）}}
\begin{center}
\begin{tikzpicture}[x=1cm,y=1cm,font=\small,>=Stealth,
  every node/.style={inner sep=3pt},line cap=round,line join=round]
\path[use as bounding box] (-1,-1.6) rectangle (11,5.5);
\node[font=\small\bfseries,text=CStatement] at (5,5.1) {以斜边 OA 为底：高必须垂直于 OA};

\coordinate (O) at (1,0); \coordinate (A) at (6,1);
\coordinate (B) at (3,3); \coordinate (H) at (3.5,0.5);
\fill[CStatement!17] (O)--(A)--(B)--cycle;
\draw[very thick,CStatement] (O)--(A);
\draw[thick] (A)--(B)--(O);
\draw[dashed,very thick,CTrap] (B)--node[right] {$d$}(H);
\pic[draw,angle radius=3mm] {right angle=A--H--B};
\node[below left] at (O) {$O$}; \node[right] at (A) {$A$};
\node[above] at (B) {$B$}; \node[below] at (H) {$H$};
\node[below] at (4,-0.25) {$|OA|=\sqrt{x_1^2+y_1^2}$};
\node[align=center] at (8.5,2.8) {$BH\perp OA$\\[8pt]$S=\frac12|OA|\cdot BH$};

\node[align=center] at (5,-1.15) {垂足 $H$ 可以在底边延长线上；距离公式始终取非负值。};
\end{tikzpicture}
\end{center}

当 $A\ne O$ 时，直线 $OA$ 为 $y_1x-x_1y=0$。
\[
 |OA|=\sqrt{x_1^2+y_1^2},\qquad
 d(B,OA)=\frac{|y_1x_2-x_1y_2|}{\sqrt{x_1^2+y_1^2}}.
\]
用 $S=\tfrac12\,|OA|\,d$，平方根约去：
\[
 S=\frac12\sqrt{x_1^2+y_1^2}\,
 \frac{|y_1x_2-x_1y_2|}{\sqrt{x_1^2+y_1^2}}
 =\boxed{\frac12|x_1y_2-x_2y_1|}.
\]
若 $A=O$，两边均为零，仍成立。
\textbf{适用信号：} 已知点到直线距离公式，不愿作辅助点。

\medskip
\end{proofbox}
\clearpage
\begin{proofbox}
\textbf{\textcolor{CProof}{方法七：平行四边形 + 正弦差角（独立推导）}}
\begin{center}
\begin{tikzpicture}[x=1cm,y=1cm,font=\small,>=Stealth,
  every node/.style={inner sep=3pt},line cap=round,line join=round]
\path[use as bounding box] (-1,-1.6) rectangle (11,5.5);
\node[font=\small\bfseries,text=CStatement] at (5,5.1) {方向角之差给出平行四边形的高};

\coordinate (O) at (1,0); \coordinate (A) at (6,0.5);
\coordinate (B) at (3,3); \coordinate (D) at (8,3.5);
\coordinate (H) at (3.277228,0.227723);
\fill[CStatement!17] (O)--(A)--(B)--cycle;
\fill[CExample!15] (A)--(D)--(B)--cycle;
\draw[thick] (O)--(A)--(D)--(B)--cycle;
\draw[CProof,thick] (A)--(B);
\draw[->,gray] (O)--(7,0) node[right] {$x$};
\draw[->,very thick,CStatement] (O)--(A);
\draw[->,very thick,CExample] (O)--node[left] {$s$}(B);
\draw[dashed,CTrap,thick] (B)--(H);
\pic[draw,angle radius=3mm] {right angle=A--H--B};
\draw[CProof] (1.995,0.0995) arc[start angle=5.7106,end angle=56.3099,radius=1];
\node at (2.4,0.65) {$\beta-\alpha$};
\node[below left] at (O) {$O$}; \node[right] at (A) {$A$};
\node[above] at (B) {$B$}; \node[right] at (D) {$D$};
\node at (3.277,-0.4) {$H$};
\node at (5,-0.45) {$r$};
\node[align=center] at (8,4.3) {$BH=s|\sin(\beta-\alpha)|$};

\node[align=center] at (5,-1.15) {$\alpha=\arg OA,\ \beta=\arg OB$；整体面积 $r\cdot BH=2S$。};
\end{tikzpicture}
\end{center}

令 $D=A+B$；$OADB$ 为平行四边形，故 $S_{OADB}=2S_{\triangle OAB}$。
若 $OA,OB$ 的方向角分别为 $\alpha,\beta$、长度为 $r,s$，
那么 $x_1=r\cos\alpha$、$y_1=r\sin\alpha$、
$x_2=s\cos\beta$、$y_2=s\sin\beta$，而
\[
\begin{aligned}
 2S_{\triangle OAB}&=S_{OADB}=rs|\sin(\beta-\alpha)|\\
 &=rs|\cos\alpha\sin\beta-\sin\alpha\cos\beta|\\
 &=|x_1y_2-y_1x_2|.
\end{aligned}
\]
故结论成立；若 $r=0$ 或 $s=0$，直接为零。
\textbf{几何意义：} 平行四边形面积 $=$ 底 $\times$ 高，差角公式把“夹角”翻译成“坐标交叉乘减”。

\medskip
\end{proofbox}
\clearpage
\begin{proofbox}
\textbf{\textcolor{CProof}{推广：任意三点的平移法}}
\begin{center}
\begin{tikzpicture}[x=1cm,y=1cm,font=\small,>=Stealth,
  every node/.style={inner sep=3pt},line cap=round,line join=round]
\path[use as bounding box] (-1,-1.6) rectangle (11,5.5);
\node[font=\small\bfseries,text=CStatement] at (5,5.1) {三个顶点同时平移，形状与面积保持不变};

\begin{scope}[x=0.65cm,y=0.65cm,shift={(-0.5,0)}]
\draw[->,gray] (0,0)--(5,0) node[right] {$x$};
\draw[->,gray] (0,0)--(0,7.3) node[above] {$y$};
\filldraw[fill=CStatement!17,draw=CStatement,thick] (1,2)--(4,5)--(2,7)--cycle;
\node[below right] at (1,2) {$P(1,2)$}; \node[right] at (4,5) {$Q(4,5)$};
\node[right] at (2,7) {$R(2,7)$};
\end{scope}
\begin{scope}[x=0.65cm,y=0.65cm,shift={(10,0)}]
\draw[->,gray] (-0.3,0)--(4,0) node[right] {$x'$};
\draw[->,gray] (0,-0.3)--(0,6) node[above] {$y'$};
\filldraw[fill=CStatement!17,draw=CStatement,thick] (0,0)--(3,3)--(1,5)--cycle;
\node[below left] at (0,0) {$O$}; \node[right] at (3,3) {$Q'(3,3)$};
\node[right] at (1,5) {$R'(1,5)$};
\end{scope}
\draw[->,thick,CProof] (4,1.7)--node[above,align=center] {同时减去\\$(1,2)$}(6,1.7);

\node[align=center] at (5,-1.15) {两图使用相同尺度；$S=\frac12|3\cdot5-1\cdot3|=6$。};
\end{tikzpicture}
\end{center}

设 $P(x_1,y_1)$、$Q(x_2,y_2)$、$R(x_3,y_3)$。
将三个顶点同时减去 $P$ 的坐标，得到
$O$、$Q'=(x_2-x_1,y_2-y_1)$、$R'=(x_3-x_1,y_3-y_1)$。
平移不改变面积，所以
\[
 \boxed{S_{\triangle PQR}=\frac12\left|
 (x_2-x_1)(y_3-y_1)-(x_3-x_1)(y_2-y_1)\right|.}
\]
\end{proofbox}
